% Lapisan solusi O001 -- Bab 14 (Aljabar Pengali)
% 2 solusi asli terpisah; bukan tulisan John M. Erdman.
% Provenans model: OpenAI Codex gpt-5.6-sol, Ultra, atas arahan pengguna.
% Lisensi komponen solusi: CC BY-SA 4.0.
% Tidak menyiratkan dukungan John M. Erdman atau Portland State University.

\clearpage
\markboth{SOLUSI O001 UNTUK BAB 14: ALJABAR PENGALI}{SOLUSI O001 UNTUK BAB 14: ALJABAR PENGALI}
% AMSBook mengambil kepala halaman genap dari topmark halaman sebelumnya.
% Gaya satu halaman ini memastikan halaman pembuka Bab 14 memakai kepala
% Bab 14 tanpa memindahkan tanda tersebut ke halaman terakhir Bab 13.
\makeatletter
\def\ps@oonechapterxiv{%
  \ps@headings
  \def\@evenhead{%
    \setTrue{runhead}%
    \normalfont\scriptsize
    \rlap{\thepage}\hfil
    \def\thanks{\protect\thanks@warning}%
    SOLUSI O001 UNTUK BAB 14: ALJABAR PENGALI\hfil}%
  \def\@oddhead{%
    \setTrue{runhead}%
    \normalfont\scriptsize\hfil
    \def\thanks{\protect\thanks@warning}%
    SOLUSI O001 UNTUK BAB 14: ALJABAR PENGALI\hfil\llap{\thepage}}%
}
\makeatother
\thispagestyle{oonechapterxiv}
\section*{Solusi O001 untuk Bab 14: Aljabar Pengali}
\addcontentsline{toc}{section}{Solusi O001 untuk Bab 14: Aljabar Pengali}
\begin{center}
\small
Komponen solusi asli terpisah, CC BY-SA 4.0. Ditulis dengan bantuan
OpenAI Codex gpt-5.6-sol, Ultra, atas arahan pengguna. Pernyataan latihan
berasal dari edisi terjemahan karya John M. Erdman; solusi ini bukan tulisan
beliau dan tidak menyiratkan dukungan beliau atau Portland State University.
\end{center}

% O001-SOLUTION-ID: O001-FAOA-2015-CH14-EX-001-SOLUTION
% SOURCE-EXERCISE-ID: FAOA-2015-CH14-NODE-0005
% STATEMENT-TARGET-SHA256: cd4d713e827bc5f3452e91f669d72db723d317f124b43c853eda067d3b95539e
\begin{o001solution}
  {O001-FAOA-2015-CH14-EX-001-SOLUTION}
  {FAOA-2015-CH14-NODE-0005}
  {cd4d713e827bc5f3452e91f669d72db723d317f124b43c853eda067d3b95539e}
\begin{o001statement}
Pastikan bahwa definisi \emph{ruang vektor} sebelumnya ekuivalen
dengan definisi yang biasa Anda gunakan.
\end{o001statement}
\begin{o001answer}
Pemetaan $M$ dan perkalian skalar biasa saling menentukan melalui
$\alpha x=M(\alpha)(x)$; aksioma homomorfisme gelanggang beridentitas bagi
$M$ tepat ekuivalen dengan aksioma perkalian skalar.
\end{o001answer}
\begin{o001proof}
Andaikan mula-mula $(V,+,M)$ memenuhi definisi formal dan tetapkan
$\alpha x=M(\alpha)(x)$. Karena setiap $M(\alpha)$ merupakan endomorfisme
grup Abel~$(V,+)$,
\[
 \alpha(x+y)=\alpha x+\alpha y.
\]
Karena $M$ homomorfisme gelanggang,
\begin{align*}
 (\alpha+\beta)x
 &=M(\alpha+\beta)x
  =\bigl(M(\alpha)+M(\beta)\bigr)x
  =\alpha x+\beta x,\\
 (\alpha\beta)x
 &=M(\alpha\beta)x
  =M(\alpha)M(\beta)x
  =\alpha(\beta x),\\
 1x&=M(1)x=x.
\end{align*}
Grup Abel memberi seluruh aksioma penjumlahan vektor, dan ketiga identitas di
atas beserta aditivitas $M(\alpha)$ memberi seluruh aksioma perkalian skalar
kompleks yang biasa.

Sebaliknya, mulai dari ruang vektor kompleks dalam pengertian biasa dan
definisikan $M(\alpha)(x)=\alpha x$. Distributivitas terhadap penjumlahan
vektor menunjukkan $M(\alpha)\in\Hom(V,+)$. Distributivitas terhadap
penjumlahan skalar memberi $M(\alpha+\beta)=M(\alpha)+M(\beta)$,
asosiativitas perkalian skalar memberi
$M(\alpha\beta)=M(\alpha)M(\beta)$, dan aksioma satuan memberi
$M(1)=I_V$. Jadi $M$ homomorfisme gelanggang beridentitas. Kedua konstruksi
jelas saling invers.
\end{o001proof}
\end{o001solution}

% O001-SOLUTION-ID: O001-FAOA-2015-CH14-EX-002-SOLUTION
% SOURCE-EXERCISE-ID: FAOA-2015-CH14-NODE-0009
% STATEMENT-TARGET-SHA256: 167b8865f270fa360a46726a6ce2259e3433b08b036fc661bf498ef6242b93dc
\begin{o001solution}
  {O001-FAOA-2015-CH14-EX-002-SOLUTION}
  {FAOA-2015-CH14-NODE-0009}
  {167b8865f270fa360a46726a6ce2259e3433b08b036fc661bf498ef6242b93dc}
\begin{o001statement}
Periksa bahwa definisi \emph{modul-$A$} dalam~\ref{0038017}
dan~\ref{0038031} ekuivalen.
\end{o001statement}
\begin{o001answer}
Dengan perkalian operator standar, kedua definisi tidak ekuivalen persis
seperti tercetak: definisi~\ref{0038031} harus memakai antihomomorfisme
$\Phi\colon A\to\ofml L(V)$, atau secara ekuivalen homomorfisme
$A^{\mathrm{op}}\to\ofml L(V)$. Setelah koreksi itu, korespondensinya ialah
$\Phi(a)x=xa$.
\end{o001answer}
\begin{o001proof}
Mulailah dengan modul kanan dalam~\ref{0038017} dan definisikan operator
$\Phi(a)$ melalui
\[
 \Phi(a)x=xa.
\]
Bilinearitas aksi menunjukkan bahwa $\Phi(a)$ linear dan bahwa
$a\mapsto\Phi(a)$ linear. Akan tetapi, aksioma modul kanan memberi
\[
 \Phi(ab)x=x(ab)=(xa)b=\Phi(b)\Phi(a)x,
\]
jadi
\[
 \Phi(ab)=\Phi(b)\Phi(a).
\]
Dengan kata lain, $\Phi$ adalah antihomomorfisme
$A\to\ofml L(V)$, atau homomorfisme
$A^{\mathrm{op}}\to\ofml L(V)$. Jika $A$ beridentitas, syarat
$x1_A=x$ tepat sama dengan $\Phi(1_A)=I_V$.

Sebaliknya, andaikan $\Phi\colon A\to\ofml L(V)$ linear, beridentitas bila
diperlukan, dan memenuhi $\Phi(ab)=\Phi(b)\Phi(a)$. Tetapkan
$xa=\Phi(a)x$. Linearitas $\Phi$ dan setiap operator $\Phi(a)$ memberi
bilinearitas, sedangkan
\[
 (xa)b=\Phi(b)\Phi(a)x=\Phi(ab)x=x(ab).
\]
Syarat satuan juga langsung cocok. Kedua konstruksi ini saling invers, sehingga
memberi ekuivalensi yang benar.

Bahwa kata ``homomorfisme'' pada~\ref{0038031} tidak dapat dipertahankan untuk
aljabar tak komutatif dapat dilihat dari contoh $A=M_2(\C)$ dan $V=A$ dengan
aksi kanan biasa. Jika $R_a(x)=xa$, maka
$R_aR_b=R_{ba}$, bukan $R_{ab}$ pada umumnya. Misalnya, untuk
$a=E_{12}$ dan $b=E_{21}$, diperoleh $ab=E_{11}$ dan $ba=E_{22}$.
Jadi pernyataan tercetak memerlukan koreksi arah perkalian di atas; solusi ini
tidak menyamarkan ketidaksesuaian tersebut.
\end{o001proof}
\end{o001solution}
